% !TEX root = ../Thesis.tex
\chapter{Abstract}

Digital collections allow people to access artworks without the space limits of a physical exhibition. However, most search systems still present their results in lists or grids. This gives a limited feeling of visiting an exhibition and makes it difficult to explore a collection from one interesting object to the next. This thesis looks at how multimedia search can be combined with Virtual Reality to present search results as dynamic exhibitions.

The concept uses one virtual museum room that can be reused for every search. Images and videos are shown on the walls, while 3D models have their own display areas. A text query creates the first exhibition. After that, visitors can select an exhibit and use it for a similarity search, which creates the next exhibition. A prototype based on this concept was developed in Godot with OpenXR and connected to the vitrivr-engine. It supports controllers and hand tracking and can be used on standalone VR headsets or through PC streaming.

A user evaluation looked at how people used the museum for the first time. It included entering the museum, searching the collection, moving through the room, and interacting with the different media types. The collected feedback gives a basis for improving the prototype and for future studies. The project provides a working example of how multimedia retrieval and spatial exploration can be brought together in VR.
